% Figure: midship section of a loaded rectangular pontoon, upright. Marks the
% beam B, the draft T, the keel K, the centre of buoyancy KB = T/2, the centre
% of gravity KG, and the metacentre M at KB + BM. The waterplane is the section
% at the waterline of width B. Used by the second end-to-end case study. The
% figure is upright (not heeled); the heeled geometry is in fig-metacentre.
% No em-dashes. Every begin matched by an end; coordinates are literal.
\begin{figure}[htbp]
\centering
\begin{tikzpicture}[scale=1.05]
% water region either side of the hull, light shade
\fill[engblue!12] (-3.6,-0.2) rectangle (3.6,1.6);
% waterline across the full width
\draw[engblue, thick] (-3.6,1.6) -- (3.6,1.6);
\node[engblue, font=\scriptsize, anchor=south east] at (3.55,1.6) {waterline};
% pontoon box: beam from -2 to 2, depth from keel at 0 to deck at 2.6,
% draft to the waterline at z = 1.6 (so T = 1.6 in figure units)
\draw[engblack, very thick] (-2.0,0.0) rectangle (2.0,2.6);
\node[engblack, font=\small] at (-1.4,2.25) {hull};
% keel point K at the base centreline
\fill[engblack] (0,0) circle (1.6pt);
\node[engblack, font=\small, anchor=north] at (0,-0.05) {$K$};
% body centreline
\draw[enggray, dashed] (0,0) -- (0,3.5);
% centre of buoyancy B at half the draft above the keel
\fill[englime] (0,0.8) circle (1.9pt);
\node[englime, font=\small, anchor=east] at (-0.08,0.8) {$B$};
% centre of gravity G above B (loaded case carries G high)
\fill[engred] (0,1.9) circle (1.9pt);
\node[engred, font=\small, anchor=west] at (0.1,1.9) {$G$};
% metacentre M above B by BM = I_wp / V_disp
\fill[engamber] (0,3.05) circle (2.0pt);
\node[engamber, font=\small, anchor=west] at (0.1,3.05) {$M$};
% dimension: beam B across the deck
\draw[enggray, <->] (-2.0,2.85) -- (2.0,2.85);
\node[enggray, font=\scriptsize, anchor=south] at (0,2.85) {beam $B$};
% dimension: draft T from keel to waterline
\draw[enggray, <->] (-2.55,0.0) -- (-2.55,1.6);
\node[enggray, font=\scriptsize, anchor=east] at (-2.55,0.8) {draft $T$};
% dimension: depth D from keel to deck
\draw[enggray, <->] (2.55,0.0) -- (2.55,2.6);
\node[enggray, font=\scriptsize, anchor=west] at (2.55,1.3) {depth $D$};
% bracket KB
\draw[englime, <->] (0.35,0.0) -- (0.35,0.8);
\node[englime, font=\scriptsize, anchor=west] at (0.4,0.4) {$\overline{KB}$};
% bracket GM (M above G, stable)
\draw[engamber, <->] (-0.45,1.9) -- (-0.45,3.05);
\node[engamber, font=\scriptsize, anchor=east] at (-0.5,2.45) {$\overline{GM}$};
\end{tikzpicture}
\caption[Midship section of a loaded pontoon]{Upright midship section of
the loaded pontoon of the second case study. The keel reference $K$ sits
on the base centreline; the centre of buoyancy $B$ (green) is at half the
draft, $\overline{KB}=T/2$; the centre of gravity $G$ (red) rides higher
because the deck load lifts it; the metacentre $M$ (amber) sits a distance
$\overline{BM}=I_{\mathrm{wp}}/V_{\mathrm{disp}}$ above $B$. The
metacentric height $\overline{GM}=\overline{KB}+\overline{BM}-\overline{KG}$
is the amber gap, positive here, so the pontoon is stable in roll. The
beam $B$ enters $\overline{BM}$ as the cube of the waterplane width, which
is why a wide flat pontoon carries large reserve stability on a shallow
draft.}
\label{fig:vol03ch10:pontoon-section}
\end{figure}
